% Add these packages to your preamble:
% \usepackage{xcolor}
% \usepackage{colortbl}
% \usepackage{longtable}


\begin{table}[htbp]
\centering
\small
\caption{Comparative Analysis: Notch Signaling Example (GoldSack Success)}
\label{tab:goldsack-success}
\begin{tabular}{|p{12cm}|}
\hline
\textbf{Success Case: Notch Signaling \& DLL Ligands (GoldSack BLEU: 20.43, YAKE BLEU: 7.02)} \\
\hline
\textbf{Reference:} A small number of signaling systems control how an animal develops from a single cell into a complex organism made up of many different cell types. Signals pass back and forth between cells, switching genes on and off to direct the development of tissues and organs. One of these signaling systems, called Notch, is so ancient that it appears in nearly all multicellular organisms. A cell sends a Notch signal using proteins called Delta or Jagged ligands that span membrane of the cell, so that part of the protein sits inside the cell and part remains outside. To change the behavior of another cell, the ligands bind to proteins called Notch receptors that span the membrane of the receiving cell. Mammals have two types of Delta ligand, two types of Jagged ligand and four types of Notch receptor... \\
\hline
\textbf{GoldSack:} \colorbox{red!20}{protein protein of cells} \colorbox{red!20}{molecules are the cells} animal develops, \colorbox{green!20}{its single cell to a complex} \colorbox{green!20}{of up of many different cells} \colorbox{green!20}{types}. \colorbox{green!20}{Onealing are between and} \colorbox{green!20}{forth between cells in and on} \colorbox{green!20}{on and off, regulate the} \colorbox{green!20}{development of tissues and} \colorbox{green!20}{organs}. \colorbox{green!20}{These of these signaling} \colorbox{green!20}{pathways is called Notch} \colorbox{green!20}{signaling is responsible} \colorbox{green!20}{important that it was to the} \colorbox{green!20}{every} animalsicellular organisms. \colorbox{green!20}{Not protein receives a} \colorbox{green!20}{signalch signal to a called} \colorbox{green!20}{receptors lig Seragged} \colorbox{green!20}{ligands}, bind the- the cell. and that the of the receptor is on the cell. interacts of outside the The activate the activity of a cell, \colorbox{green!20}{the ligands bind to the on} \colorbox{green!20}{Dch receptors on are the} \colorbox{green!20}{membrane of the cell cell}. Theals have two Not of Not ligands, D of of Jagged ligand and two types of Notch lig... \\
\hline
\textbf{YAKE:} \colorbox{red!20}{protein signal of cells} \colorbox{red!20}{molecules in the cells} animal develops, \colorbox{red!20}{its single fertil to a complex} \colorbox{red!20}{of of of many different} \colorbox{red!20}{types types}. \colorbox{red!20}{Oneals from from and forth} \colorbox{red!20}{between the as and on on} \colorbox{red!20}{and off}, \colorbox{red!20}{control the formation of} \colorbox{red!20}{different and organs}. \colorbox{red!20}{The of the signaling} \colorbox{red!20}{pathways is called Notch} \colorbox{red!20}{signaling is involved} \colorbox{red!20}{important that it is in the} \colorbox{red!20}{all} animalsicellular organisms. \colorbox{red!20}{The receptor detects a} \colorbox{red!20}{signalch signal to a on lig-} \colorbox{red!20}{Lagagged ligands}, bind the on the receptor. and they they of the receptor on on the membrane membrane allows of outside. \colorbox{red!20}{The bind the activity of the} \colorbox{red!20}{lig, the ligands must to the} \colorbox{red!20}{on Notch receptors on are} \colorbox{red!20}{the surface of the receptor} \colorbox{red!20}{cell}. Theals have two Not of lig ligands, D that of Jagged ligand and one activating of Lagch receptor... \\
\hline
\end{tabular}
\end{table}

\begin{table}[htbp]
\centering
\small
\caption{Comparative Analysis: Neuron Synapses Example (YAKE Success)}
\label{tab:yake-success}
\begin{tabular}{|p{12cm}|}
\hline
\textbf{Success Case: Neuron Synapses (YAKE BLEU: 15.91, GoldSack BLEU: 8.80)} \\
\hline
\textbf{Reference:} Neurons communicate with one another at junctions called synapses. When an electrical signal travels along a neuron and arrives at a synapse, vesicles filled with small neurotransmitter molecules fuse with the cell membrane and release the neurotransmitter. These chemicals rapidly bind to receptors on the downstream neuron that induce an electrical response in that cell. Vesicles can be consumed at prodigious rates, up to 500 a second, so the cell must recover the membrane rapidly and regenerate more vesicles filled with neurotransmitter. Experiments in the 1970s and 1980s suggested that when vesicles empty their contents into the synapse, they fuse completely with the membrane and are lost. To recover the membrane, the cell forms 'pits', by means of a coat protein called clathrin... \\
\hline
\textbf{GoldSack:} \colorbox{red!20}{urons communicate with} \colorbox{red!20}{each another by} thections called synapses. \colorbox{red!20}{The a electrical signal is} \colorbox{red!20}{between a nerve, reaches} \colorbox{red!20}{at a newapse, itesicles} \colorbox{red!20}{called with a moleculesmit} \colorbox{red!20}{molecules} are with the membrane membrane and release \colorbox{red!20}{the signalmitter}. \colorbox{red!20}{The v are travel to each on} \colorbox{red!20}{the membrane cell} and \colorbox{red!20}{are the electrical response.} \colorbox{red!20}{the neuron}. \colorbox{red!20}{Theicles are be released by} \colorbox{red!20}{synigious rates}, which to 100 times second, and it recycling must be them v- to efficiently the vesicles. with neurotransmit. Howeveriments in frog 1970s suggested 1980s suggested that v neuronsesicles are from v at the plasmaapse, they are with with the cell and release recycled. However test the membrane, v v must a kissockets' that which which of which protein of called clathrin, which coats coatss from a membrane... \\
\hline
\textbf{YAKE:} \colorbox{green!20}{urons communicate with one} \colorbox{green!20}{another at junctions called} \colorbox{green!20}{synapses}. \colorbox{green!20}{When a electrical impulse} \colorbox{green!20}{reaches from a neuron,} \colorbox{green!20}{reaches at a} secondapse, \colorbox{green!20}{itesicles containing with} bubble bubblemitter molecules \colorbox{green!20}{fuse with the membrane} membrane and \colorbox{green!20}{release the} contentsmitter. \colorbox{green!20}{These v are leave to} \colorbox{green!20}{receptors on the surface} side and relay movement electrical response. the neuron. \colorbox{green!20}{Theicles are be released by} \colorbox{green!20}{synapses rates} and which to 100 times second, and it neurons needs be them v-. recycle the vesicles. with neurotransmitter. Howeveriments in frogs laboratorys and 1980s suggested that v neuronsesicles are at membrane at the membraneapses, they are with with the membrane. are recycled. However test the membrane, v v must a kissou' that which which of which protein of called clathrin. which is bindss and the membrane... \\
\hline
\end{tabular}
\end{table}

% Key observations:
% - Green highlighting: Better phrasing compared to the other approach
% - Red highlighting: Errors in UMLS entity resolution or grammatical issues
% 
% Example 110 (GoldSack wins - BLEU 20.43 vs 7.02): GoldSack maintains better core biological 
% terminology and sentence flow like "single cell to a complex...many different cells types", 
% "regulate the development of tissues and organs", "Notch signaling". YAKE corrupts heavily 
% with "fertil to a complex", "Oneals from from", "Lajagged ligands"
%
% Example 164 (YAKE wins - BLEU 15.91 vs 8.80): YAKE preserves better sentence structure like 
% "Neurons communicate with one another at junctions called synapses" and "vesicles...fuse with 
% the membrane", while GoldSack fragments heavily with "newapse, itesicles" and "signalmitter"
